%package list
\documentclass{article}
\usepackage[top=3cm, bottom=3cm, outer=3cm, inner=3cm]{geometry}
\usepackage{graphicx}
\usepackage{url}
%\usepackage{cite}
\usepackage{hyperref}
\usepackage{array}
\usepackage{amsmath}
%\usepackage{multicol}
\newcolumntype{x}[1]{>{\centering\arraybackslash\hspace{0pt}}p{#1}}
\usepackage{natbib}
\usepackage{pdfpages}
\usepackage{multirow}
\usepackage{multirow}
\usepackage[normalem]{ulem}
\useunder{\uline}{\ul}{}


%%%%%%%%%%%%%%%%%%%%%%%%%%%%%%%%%%%%%%%%%%%%%%%%%%%%%%%%%%%%%%%%%%%%%%%%%%%%
%%%%%%%%%%%%%%%%%%%%%%%%%%%%%%%%%%%%%%%%%%%%%%%%%%%%%%%%%%%%%%%%%%%%%%%%%%%%
\newcommand{\csemail}{vmachacaa@ulasalle.edu.pe}
\newcommand{\csdocente}{MSc. Vicente Enrique Machaca Arceda}
\newcommand{\cscurso}{Fundamentos de Lenguajes de Programación}
\newcommand{\csuniversidad}{Universidad La Salle}
\newcommand{\csescuela}{Escuela Profesional de Ingeniería de Software}
\newcommand{\cspracnr}{05}
\newcommand{\cstema}{Ensamblador}
%%%%%%%%%%%%%%%%%%%%%%%%%%%%%%%%%%%%%%%%%%%%%%%%%%%%%%%%%%%%%%%%%%%%%%%%%%%%
%%%%%%%%%%%%%%%%%%%%%%%%%%%%%%%%%%%%%%%%%%%%%%%%%%%%%%%%%%%%%%%%%%%%%%%%%%%%


\usepackage[english,spanish]{babel}
\usepackage[utf8]{inputenc}
\AtBeginDocument{\selectlanguage{spanish}}
\renewcommand{\figurename}{Figura}
\renewcommand{\refname}{Referencias}
\renewcommand{\tablename}{Tabla} %esto no funciona cuando se usa babel
\AtBeginDocument{%
	\renewcommand\tablename{Tabla}
}

\usepackage{fancyhdr}
\pagestyle{fancy}
\fancyhf{}
\setlength{\headheight}{30pt}
\renewcommand{\headrulewidth}{1pt}
\renewcommand{\footrulewidth}{1pt}
\fancyhead[L]{\raisebox{-0.2\height}{\includegraphics[width=3cm]{img/logo_salle}}}
\fancyhead[C]{}
\fancyhead[R]{\fontsize{7}{7}\selectfont	\csuniversidad \\ \csescuela \\ \textbf{\cscurso} }
\fancyfoot[L]{MSc. Vicente Machaca}
\fancyfoot[C]{FLP}
\fancyfoot[R]{Página \thepage}


\usepackage{listings} 
\usepackage{color}
\definecolor{dkgreen}{rgb}{0,0.6,0}
\definecolor{gray}{rgb}{0.5,0.5,0.5}
\definecolor{mauve}{rgb}{0.58,0,0.82}
\lstset{frame=tb,
	language=Python,
	aboveskip=3mm,
	belowskip=3mm,
	showstringspaces=false,
	columns=flexible,
	basicstyle={\small\ttfamily},
	numbers=none,
	numberstyle=\tiny\color{gray},
	keywordstyle=\color{blue},
	commentstyle=\color{dkgreen},
	stringstyle=\color{mauve},
	breaklines=true,
	breakatwhitespace=true,
	tabsize=3
}




\begin{document}
	
	\vspace*{10px}
	
	\begin{center}	
		\fontsize{17}{17} \textbf{ Práctica \cspracnr}
	\end{center}
	%\centerline{\textbf{\underline{\Large Título: Informe de revisión del estado del arte}}}
	%\vspace*{0.5cm}
	

	\begin{table}[h]
		\begin{tabular}{|x{4.7cm}|x{4.8cm}|x{4.8cm}|}
			\hline 
			\textbf{DOCENTE} & \textbf{CARRERA}  & \textbf{CURSO}   \\
			\hline 
			\csdocente & \csescuela & \cscurso    \\
			\hline 
		\end{tabular}
	\end{table}	
	
	
	\begin{table}[h]
		\begin{tabular}{|x{4.7cm}|x{4.8cm}|x{4.8cm}|}
			\hline 
			\textbf{PRÁCTICA} & \textbf{TEMA}  & \textbf{DURACIÓN}   \\
			\hline 
			\cspracnr & \cstema & 3 horas   \\
			\hline 
		\end{tabular}
	\end{table}
	
	
	\section{Competencias del curso}
	\begin{itemize}
		\item Conocer el desarrollo histórico de los lenguajes de programación y los paradigmas de programación.
		\item Comprender el papel de los diferentes mecanismos de abstracción en la creación de facilidades definidas por el usuario así como los beneficios de los lenguajes intermedios en el proceso de compilación.		
	\end{itemize}
	
	
	\section{Competencias de la práctica}
	\begin{itemize}
		\item Programar tareas básicas en lenguaje ensamblador.
	\end{itemize}
	
	\section{Equipos y materiales}
	\begin{itemize}
		\item Latex
		\item Conección a internet 
		\item IDE de desarrollo
		%\item Matplotlib 
		%\item Numpy 
		%\item BioPython
		%\item Cuenta en Github
	\end{itemize}

	\section{Entregables}
	\begin{itemize}		
		\item Se debe elaborar un informe en \textbf{Latex} donde se responda a cada ejercicio de la Sección \ref{sec:ejercicios}.
		\item En el informe se debe agregar un enlace al repositorio Github donde esta el código.
		\item En el informe se debe agregar el código fuente asi como capturas de pantalla de la ejecución y resultados del mismo.
		\item Por cada 5 minutos de retraso, el alumno tendrá un punto menos.
	\end{itemize}


	
	\clearpage
	
	\section{Ejercicios}\label{sec:ejercicios}	
	
	\begin{enumerate}
		\item Escribir un programa que pida al usuario un número entero y muestre por pantalla si es par o impar. \textbf{(3 puntos)}
		
		\begin{lstlisting}
		>> Ingrese un numero: 3
		>> El numero 3 es impar		
		\end{lstlisting}
		
		\begin{lstlisting}
		>> Ingrese un numero: 4
		>> El numero 4 es par		
		\end{lstlisting}
		
		\item Escribir un programa que pida al usuario un número entero positivo y muestre por pantalla todos los números impares desde 1 hasta ese número. \textbf{(4 puntos)}
		
		\begin{lstlisting}
		>> Ingrese un numero: 10
		>> Los numeros impares hasta 10 son: 1 3 5 7 9	
		\end{lstlisting}
		
		\item Escribir un programa que pida al usuario un número entero y muestre por pantalla si es un número primo o no. \textbf{(5 puntos)}
		
		\begin{lstlisting}
		>> Ingrese un numero: 10
		>> El numero 10 no es primo	
		\end{lstlisting}
		
		\begin{lstlisting}
		>> Ingrese un numero: 7
		>> El numero 7 es primo	
		\end{lstlisting}
		
		\item Escriba un programa que entrege el máximo común divisor de dos enteros ingresados por teclado. \textbf{(5 puntos)}
				
		\begin{lstlisting}
		>> Ingrese un numero: 5
		>> Ingrese un numero: 15		
		>> El maximo comun divisor es: 5
		\end{lstlisting}
		
		\begin{lstlisting}
		>> Ingrese un numero: 8
		>> Ingrese un numero: 12		
		>> El maximo comun divisor es: 4
		\end{lstlisting}

	\end{enumerate}


	
\clearpage
	\section{Rúbricas}
	
	\begin{table}[hbt!]
		\setlength{\tabcolsep}{0.5em} % for the horizontal padding
		{\renewcommand{\arraystretch}{1.5}% for the vertical padding
		\begin{tabular}{|p{5cm}|x{3cm}|x{3cm}|x{3cm}|}
			\hline 
			\textbf{Rúbrica} & \textbf{Cumple}  & \textbf{Cumple con obs.}  & \textbf{No cumple} \\
			\hline 
			\textbf{Informe}: El informe debe estar en Latex, con un formato limpio y facil de leer.  & 3 & 1.5 & 0   \\ 
			\hline 
			\textbf{Implementación}: Implementa la funcionalidad correcta de cada ejercicio.  & 12 & 6 & 0   \\ \hline
			\textbf{Mensaje de ayuda}: Agrega mensajes de ayuda (textos solicitando que se ingrese los números, etc.) por consola.  & 5 & 2.5 & 0   \\ \hline
			
			%\textbf{Implementación}: Responde cada ejercicio correctamente con una complejidad alta.  & 10 & 5 & 0   \\ \hline
							
			%\textbf{Presentación}: El alumno demuestra dominio del tema y conoce con exactitud cada parte de su código. & 2 & 1 & 0   \\ 			\hline 
			
			\textbf{Errores ortográficos}: Por cada error ortográfico, se le descontara 1 punto.  & - & - & -   \\ \hline
			
		\end{tabular}
	}
	\end{table}
	
	
	%\bibliographystyle{apalike}
	%\bibliographystyle{IEEEtranN}
	%\bibliography{bibliography}
	
	
	
	
	
\end{document}